\documentclass[conference]{IEEEtran}
\IEEEoverridecommandlockouts
% The preceding line is only needed to identify funding in the first footnote. If that is unneeded, please comment it out.
\usepackage{cite}
\usepackage{amsmath,amssymb,amsfonts}
\usepackage{algorithmic}
\usepackage{url}
\usepackage{siunitx}
\usepackage{graphicx}
\usepackage{textcomp}
\usepackage{xcolor}


\def\BibTeX{{\rm B\kern-.05em{\sc i\kern-.025em b}\kern-.08em
    T\kern-.1667em\lower.7ex\hbox{E}\kern-.125emX}}
\begin{document}

\title{
    Design and Implementation of a High-Power Motor Controller for Bicycles
}

%\maketitle


\author{\IEEEauthorblockN{Hugo Abescat}
\IEEEauthorblockA{\textit{GEI Department} \\
\textit{INSA Toulouse}\\
Toulouse, France \\
abescat@insa-toulouse.fr}
\and
\IEEEauthorblockN{Karima Attar}
\IEEEauthorblockA{\textit{GEI Department} \\
\textit{INSA Toulouse}\\
Toulouse, France \\
karima.attar@insa-toulouse.fr}
\and
\IEEEauthorblockN{Brage Flønæs Johnsen}
\IEEEauthorblockA{\textit{GEI Department} \\
\textit{INSA Toulouse}\\
Toulouse, France \\
johnse@insa-toulouse.fr}
\and
\IEEEauthorblockN{Oskar Orvik}
\IEEEauthorblockA{\textit{GEI Department} \\
\textit{INSA Toulouse}\\
Toulouse, France \\
orvik@insa-toulouse.fr}
\and
\IEEEauthorblockN{Julien Pavillon}
\IEEEauthorblockA{\textit{GEI Department} \\
\textit{INSA Toulouse}\\
Toulouse, France \\
pavillon@insa-toulouse.fr}
\and
\IEEEauthorblockN{Nolan Reynier Nomer}
\IEEEauthorblockA{\textit{GEI Department} \\
\textit{INSA Toulouse}\\
Toulouse, France \\
reynier-nome@insa-toulouse.fr}
\and
\IEEEauthorblockN{Aleksander Taban}
\IEEEauthorblockA{\textit{GEI Department} \\
\textit{INSA Toulouse}\\
Toulouse, France \\
taban@insa-toulouse.fr}
}

\maketitle

\begin{abstract}
Electric bikes are becoming an increasingly attractive solution for transporting goods between short distances, 
especially in city-wide infrastructures. However, most commercially available controllers rely on complex integrated 
circuits making repair and local manufacturing difficult, particularly for organisations operating in 
resource-constrained or low-tech environments. La Manufacture Autonome Décentralisée (LaMAD) is developing products 
and solutions, particularly e-bikes, which are more repairable and sustainable. Previous studies have predominantly 
focused on performance optimisation of Field Oriented Control (FOC) and trapezoidal commutation strategies, with limited 
attention to repairability, component sourcing, and community-centred sustainability criteria. This project aims to 
design, assemble, and develop a functional, low-tech and open-source motor controller for electric cargo bikes. The 
current model uses an open-source motor control called VESC (Vedder Electronic Speed Controller) that allows precise 
control of electric motors. The controller needs to be compatible with a VESC controller and easily locally repairable 
by LaMAD. By exploring the inner workings of the VESC project, modelling of the physical systems and the Printed Circuit
 Board, PCB, we investigated the ways we could do it in another way. We acquired a VESC controller to compare our system 
 and a commercial product. Preliminary results demonstrate that the adapted VESC-based controller successfully drives 
 the target motor under both commutation strategies, and that positional control is achievable with the current hardware
  configuration. Security vulnerabilities related to open Bluetooth access were identified. These findings suggest that 
  open-source, locally fabricated motor controllers can meet the functional requirements of electric cargo bikes while 
  significantly improving repairability. 
\end{abstract}

\begin{IEEEkeywords}
VESC Project, Brushless DC motor, Field Oriented Control, Trapezoidal commutation, Low-Tech, e-bike.
\end{IEEEkeywords}

\section{Introduction}
The fast urbanization of global logistics has positioned electric cargo bikes as a primary solution. At the heart of 
these vehicles is the motor controller. Current research and industry standards primarily focus on two methods of 
commutation for the controller: Trapezoidal commutation and Field Oriented Control (FOC). 

As motor controllers become smarter, they increasingly incorporate wireless connectivity for tuning and diagnostics. 
Current research highlights that while Bluetooth Low Energy (BLE) and mobile app integration improve user experience, 
they often introduce vulnerabilities. Open-source projects, in particular, must balance ease of access for community 
developers with the need to secure the vehicle.

We also argue the need for general public's safety when it comes to these bikes, as it could be a danger to the traffic.
 This is especially true when it comes to vehicules carrying a substantial load. This needs to be considered by laMAD, 
 where their responsibility and control begins and ends. Should there be a difference between the firmware loaded on a 
 product from laMAD than what is publicly available? 


\section{Related Work}

\subsection{Modeling of BLDC Motor}
The electromechanical model of a BLDC motor is foundational for understanding its behavior under different control 
schemes. BLDC motors are categorized by their back-electromotive force (back-EMF) waveform: trapezoidal or sinusoidal. 
This distinction is crucial, as the trapezoidal shape inherently leads to torque ripple when the supplied phase currents 
are not perfectly aligned, directly influencing the choice and effectiveness of the control strategy 
\cite{patil_analysis_2025}. For a BLDC motor with trapezoidal back-EMF, the electromagnetic torque is given by:
\[
T_e = \frac{e_a i_a + e_b i_b + e_c i_c}{\omega_m}
\]
where \( e_x \) is the back-EMF and \( i_x \) is the phase current \cite{li_quantitative_2019}. The classical d-q 
reference frame model, ideal for sinusoidal machines, is less suitable for trapezoidal BLDC motors because it assumes 
sinusoidal flux distribution. Phase-variable modeling in the natural (abc) frame is therefore more appropriate, as it 
directly accounts for the non-sinusoidal, trapezoidal nature of the back-EMF and the associated harmonics 
\cite{mohammd_taher_new_2021}.

\subsection{Trapezoidal Commutation (Six-Step Control) for BLDC Motors}
% Trapezoidal commutation, or six-step control, uses Hall-effect sensors to synchronize phase current switching every 
% 120 electrical degrees. 
Trapezoidal commutation, or Six-Step control, uses bipolar conduction, with two motor phases conducting at any time and 
current commutation occurring every 120 electrical degrees \cite{gieras_modern_2023}. As commutation depends on rotor 
position, Six-Step control requires either position sensors (e.g. Hall sensors, encoders, or resolvers) or sensorless 
estimation based on back-EMF detection or observers \cite{gieras_modern_2023, gasc_conception_2004}.  This method is 
renowned for its simplicity of implementation and low hardware cost \cite{bhatiya_bldc_2024}. It enables effective 
torque control but introduces significant torque ripple during commutation events, especially under high load 
\cite{jomsa-nga_torque_2024}. This ripple generates noise, increases mechanical stress, and reduces overall efficiency 
\cite{mohammd_taher_new_2021}. Although PWM techniques can mitigate this ripple, they do not completely eliminate it 
\cite{li_quantitative_2019}.

\subsection{Field-Oriented Control (FOC) for BLDC Motors}
FOC is a vector control strategy that decouples the stator flux and torque components. It transforms three-phase 
currents into orthogonal \( I_d \) and \( I_q \) components, enabling precise torque control and significant ripple 
reduction \cite{jomsa-nga_torque_2024}. FOC is particularly effective for BLDC motors with sinusoidal back-EMF but can 
also be applied to trapezoidal back-EMF motors, albeit with less impressive ripple suppression results 
\cite{li_quantitative_2019}. It requires greater computational power and more precise position sensors (e.g. encoders). 
Comparative analysis shows that FOC yields a more stable stator current profile and significantly reduces torque 
variations compared to trapezoidal control \cite{patil_analysis_2025}.

\subsection{Comparative Analysis: FOC vs. Trapezoidal for Light Electric Vehicles}

\subsubsection{Torque Ripple and User Comfort}
Firstly, torque ripple can be reduced for both control methods by selecting appropriate motor parameters, such as the 
number of stator slots and rotor poles \cite{gasc_conception_2004}.
FOC substantially reduces torque ripple compared to Six-Step control, directly enhancing ride comfort and minimizing 
vibrations. Experimental results show a torque ripple of \SI{18.38}{\percent} for FOC versus \SI{35.67}{\percent} for 
Six-Step control at 500~rpm \cite{jomsa-nga_torque_2024}. 
Commutation torque ripple (CTR), prominent in Six-Step control, can be specifically targeted and mitigated using 
advanced control techniques like Model Predictive Control (MPC) while retaining the fundamental simplicity of 
trapezoidal commutation \cite{mohammd_taher_new_2021}.

\subsubsection{Energy Efficiency}
FOC optimizes torque per ampere (MTPA), improving efficiency at low loads. Six-Step control exhibits lower switching 
losses at high speeds \cite{li_quantitative_2019}.

\subsubsection{Complexity, Cost, and Low-Tech Suitability}

Six-Step control is inherently simpler, cheaper, and more robust, making it a prime candidate for low-tech applications. 
Research focused on reducing propulsion system costs proposes simplified hardware topologies, such as 4-switch inverters 
(instead of 6) coupled with direct current control strategies, maintaining acceptable performance while significantly 
lowering hardware costs \cite{lee_advanced_2001}. FOC, while superior in performance, is more complex to implement and 
carries higher hardware costs (sensors, processing power).

\subsubsection{Dynamic Response}

FOC provides faster response times and better load disturbance rejection \cite{jomsa-nga_torque_2024}.


\section{Research gap}
Despite this progress, limited research has examined the adaptation of open-source motor controllers to LowTech and 
repairability constraints. To date, researchers have not addressed the challenge of designing a controller that can 
be locally fabricated, repaired with standard components, and secured against unauthorised wireless access requirements 
that are critical for decentralised, community-operated fleets.

\section{The aim of the study}
This report presents the design and implementation of a VESC-based motor controller tailored to the needs of the 
Manufacture Autonome Décentralisée (MAD), an organisation operating electric cargo bikes and freight tricycles. 
We aim to focus on adapting the VESC open-source firmware to support both FOC and trapezoidal commutation, integrating 
positional control, and addressing Bluetooth security, while prioritising local manufacturability at INSA Toulouse.

\section{Software and Connectivity}

\subsection{BLE Compatibility With the VESC}

\subsubsection{First Experiment}
VESC-controllers are not necessarily equipped with Bluetooth-modules by default. Often, it is necessary to add a 
BLE-module. A standard HC-05 bluetooth-module compatible with arduino is a great way to send and recieve 
bluetooth-packets from a host, e.g. a mobile phone, via a bridge translating the bluetooth packets to the UART protocol. 
This could be demonstrated using a ESP8622's standard library with said module, by letting us send characters from one 
device to another.

\subsubsection{HC-05 and the VESC}
By flashing the VESC firmware on a discovery-card and connecting the HC-05 module to the PB10 and PB11-pins, which are 
the Rx and Tx-pins for the STM32F4xx chip, we discovered that the setup for the bluetooth module was not available in 
the VESC tool. The inherent BLE capabilities is an important limitation to consider when designing a VESC system.
We learned therefore that the HC-05 is not originally adapted for BLE. The need for a bridge also adds on complexity 
and cost, in the form of extra components and another device to maintain the code of. For the future, choosing a 
bluetooth module supporting BLE will be the easiest solution. Preferably a module fitting the communication connector on
 the cheap FOCer project\cite{b1} could facilitate the relevancy of the PCB project with a microcontroller.

\subsubsection{BLE Vulnerability}
Bluetooth could be a vulnerability to a VESC if it is to be used as a controller in real-time, as the controller could 
be jammed. Our test with the Flipper Zero shows the disfunctionnality of Bluetooth with different use cases. 
We experienced with the jamming of a bluetooth speaker that the music completely stopped. It could also be investigated 
how the connection to the VESC could be modified using the vesc tool. We will touch more on the accessability of the 
code within the vesc tool sooner. 

\subsection{Code integrity}
\subsubsection{Context}
As the project is open source, and the code is freely accessible, there should be no reason to hide the code. It could 
however be reasonable to protect the code from changes which could hurt other people. Changing following parameters 
should at least come with a disclaimer and clearly state the dangers possible by proceeding with said changes. We 
have in mind the maximum speed permitted and the power available to the motors.

\subsubsection{LispBM extraction}
We caugth word that the lisp code for the VESC used by Maillon mobility was easy to extract. By building an older 
firmware with the Maillon mobility software, we observed this by going to the lispBM tab and clicking read. It's up to 
laMAD if they would like to reinforce this mechanism. A modification on a parameter and then clicking upload allowed 
us to easily change the speed limit. This could bring up a public danger. This raises questions on the use of laMADs 
equipment which is in a traffic friendly manner.

\subsubsection{LispBM Code}
When we flashed newer firmware from the project made by Benjamin Vedder\cite{b1}, we also observed some difficulties in 
uploading the lispBM script taken from the one on firmware version 6.06. This could indicate that there needs to be 
further maintenance of the code in order to get the software up to speed. This needs to be documented better for someone
 to continue the project. This could be a good investment for laMAD as well in the context of training for the people 
 working on the motor control part of the e-bike.

This documentation could be as simple as referencing the relevant parts of the lispBM documentation \cite{b2}

\subsubsection{Proposed Solution}
This risk could be patched by developing a VESC application for the VESC controller or using a binary. This is a 
solution which is less open source, but which is make unlawful use of the material harder. The application could be 
created using C and use an algorithm known by laMAD in order to securise the access to someone to change the parameters 
only if they are laMAD certified personnel. This encryption would preferably be reduced to the most essential settings 
in order to align with what our impression of the philosophy of laMAD would be.



\subsection{VESC Compiling}
As mentionned, we have been able to compile the VESC tool and the VESC firmware. This firmware has been put onto an 
STM32F4xx Discovery card. This poses several obstacles for our progress on the topic of cybersecurity. We will however 
summarise what we have learned for you and propose some additional work for the future. The challenges we encountered 
were the following: The lack of bluetooth capabilities. We did not have a module with BLE either. We had access to a 
HC-05 module, but that only allows for a normal bluetooth protocol and would require further work on a bridge to UART 
by using an esp8622 that we had as well. We propose that the next group has access to a VESC controller from the 
beginning, as well as a motor we could control. This could be in cooperation with laMAD, as laMAD could propose 
some models they're interested in.

We also found that the information on the VESC is scattered around the internet. The ressources is also sometimes 
based on a debian-based linux system which adds more work for someone using another distribution of linux. This could 
hinder the implementation facility for new users. We struggeled particularly with the Qt packages for positioning and 
gamepad. We would therefore recommend the use of a debian-based linux system for the computer working with the VESC 
for the laMAD associates.


\section{Discussion}
This project could be seen as an introduction to the VESC project for someone who don't know about it from beforehand, 
the challenges the new users face during setup, as well as a demand for clear expectations concerning documentation 
on the subject. The project laMAD is leading should probably not be a fork of the project, as the project is still 
in development.

As a final note, this proved to be a project which could easily be developed into several different projects in 
different fields. Some projects could be continued later on as a different PIR subject, other could be proposed to 
later years in different spesialisations like TLS SEC, ESPE. Our thoughts on the following projects that could be 
explored are the following.

The fabrication line for electronics is globalised. This is okay in a stable world, but it could be a problem in a world 
full of instability, be it war, blockages, or tarifs. The idea of opening a spesialisation in cooperation with AIME came 
up as an idea.

For TLS SEC the subject could be the design for a fitting mechanism to restrict certain priveligies to certified 
personnel that could be used in the C programming language. Later down the line we could also see the possibility to 
analyse the Bluetooth frames in order to manipulate them in order to change important parameters.

The continuation on the PCB could be a subject fitting an ESPE spesialisation.

The proposition of and supply of a vesc system to play with and troubleshoot could be a good rule of thumb, which 
allows for a quicker start and gives among other things an idea of the budget and the supply line used by a entity 
in the sector. Proposing a visit could also be one way to familiarise students with the association.

What should be a clear conclusion from our test with the jammer is that a controlles based on Bluetooth alone should be 
avoided when possible and practical. Examples where this could be relevant include electric skateboards, as cables could 
impose a tripping hazard. There, an encapsulation of an encrypted control frame could be an thought. 

\section{Results}

\section{Conclusion/Summary}

%\begin{figure}[htbp]
%%\centerline{\includegraphics{fig1.png}}
%caption{Example of a figure caption.}
%\label{fig}
%\end{figure}

%Figure Labels: Use 8 point Times New Roman for Figure labels. Use words 
%rather than symbols or abbreviations when writing Figure axis labels to 
%avoid confusing the reader. As an example, write the quantity 
%``Magnetization'', or ``Magnetization, M'', not just ``M''. If including 
%units in the label, present them within parentheses. Do not label axes only 
%with units. In the example, write ``Magnetization (A/m)'' or ``Magnetization 
%\{A[m(1)]\}'', not just ``A/m''. Do not label axes with a ratio of 
%quantities and units. For example, write ``Temperature (K)'', not 
%``Temperature/K''.

\section*{Acknowledgment}

The preferred spelling of the word ``acknowledgment'' in America is without 
an ``e'' after the ``g''. Avoid the stilted expression ``one of us (R. B. 
G.) thanks $\ldots$''. Instead, try ``R. B. G. thanks$\ldots$''. Put sponsor 
acknowledgments in the unnumbered footnote on the first page.

%\section*{References}

%Please number citations consecutively within brackets \cite{b1}. The 
%sentence punctuation follows the bracket \cite{b2}. Refer simply to the reference 
%number, as in \cite{b3}---do not use ``Ref. \cite{b3}'' or ``reference \cite{b3}'' except at 
%the beginning of a sentence: ``Reference \cite{b3} was the first $\ldots$''

\bibliographystyle{IEEEtran}
\bibliography{PIR_MadMax3}


\end{document}
